%2multibyte Version: 5.50.0.2953 CodePage: 1252

\documentclass{article}
%%%%%%%%%%%%%%%%%%%%%%%%%%%%%%%%%%%%%%%%%%%%%%%%%%%%%%%%%%%%%%%%%%%%%%%%%%%%%%%%%%%%%%%%%%%%%%%%%%%%%%%%%%%%%%%%%%%%%%%%%%%%%%%%%%%%%%%%%%%%%%%%%%%%%%%%%%%%%%%%%%%%%%%%%%%%%%%%%%%%%%%%%%%%%%%%%%%%%%%%%%%%%%%%%%%%%%%%%%%%%%%%%%%%%%%%%%%%%%%%%%%%%%%%%%%%
\usepackage{amsfonts}
\usepackage{amsmath}

\setcounter{MaxMatrixCols}{10}
%TCIDATA{OutputFilter=LATEX.DLL}
%TCIDATA{Version=5.50.0.2953}
%TCIDATA{Codepage=1252}
%TCIDATA{<META NAME="SaveForMode" CONTENT="1">}
%TCIDATA{BibliographyScheme=Manual}
%TCIDATA{Created=Tuesday, September 13, 2016 11:26:35}
%TCIDATA{LastRevised=Monday, December 19, 2022 11:25:46}
%TCIDATA{<META NAME="GraphicsSave" CONTENT="32">}
%TCIDATA{<META NAME="DocumentShell" CONTENT="Standard LaTeX\Blank - Standard LaTeX Article">}
%TCIDATA{Language=American English}
%TCIDATA{CSTFile=40 LaTeX article.cst}

\newtheorem{theorem}{Theorem}
\newtheorem{acknowledgement}[theorem]{Acknowledgement}
\newtheorem{algorithm}[theorem]{Algorithm}
\newtheorem{axiom}[theorem]{Axiom}
\newtheorem{case}[theorem]{Case}
\newtheorem{claim}[theorem]{Claim}
\newtheorem{conclusion}[theorem]{Conclusion}
\newtheorem{condition}[theorem]{Condition}
\newtheorem{conjecture}[theorem]{Conjecture}
\newtheorem{corollary}[theorem]{Corollary}
\newtheorem{criterion}[theorem]{Criterion}
\newtheorem{definition}[theorem]{Definition}
\newtheorem{example}[theorem]{Example}
\newtheorem{exercise}[theorem]{Exercise}
\newtheorem{lemma}[theorem]{Lemma}
\newtheorem{notation}[theorem]{Notation}
\newtheorem{problem}[theorem]{Problem}
\newtheorem{proposition}[theorem]{Proposition}
\newtheorem{remark}[theorem]{Remark}
\newtheorem{solution}[theorem]{Solution}
\newtheorem{summary}[theorem]{Summary}
\newenvironment{proof}[1][Proof]{\noindent\textbf{#1.} }{\ \rule{0.5em}{0.5em}}
\input{tcilatex}
\begin{document}


\begin{center}
{\Large Exam}\bigskip 
\end{center}

\textbf{Question 1}

(a) State (without proof) the Gauss Markov and Cram\'{e}r Rao theorems, and
explain the differences between the two.

(b) State the Central Limit Theorem for a sequence of i.i.d. random
variables carefully stating any conditions required. Briefly explain the
theorem's importance for statistical inference.

(c) Suppose $\sqrt{n}\left( \hat{\theta}_{n}-\theta \right) \approx \mathcal{%
N}\left( 0,2\right) $ and you want to test a null hypothesis $H_{0}:\theta =0
$ vesus a two-sided alternative $H_{1}:\theta \neq 0.$ Explain how to
conduct a $z$ test at some significance level $1-\alpha $, and how the $z$
statistic behaves under the null and alternative.

(d) Define Type I and Type II errors in hypothesis testing and explain how
the probability of making such errors relates to the size and power of the
test.

\bigskip

\textbf{Question 2}

Consider the linear regression model $y=X\beta +\varepsilon ,$ where $y$ is
an $n\times 1$ vector of dependent variables, $X$ is an $n\times k$ matrix
of regressors, $\beta $ is a $k\times 1$ vector of unknown parameters, and $%
\varepsilon $ is an $n\times 1$ vector of unobserved random variables. The
standard assumptions provided in the lecture are assumed to hold.

(a) Suppose there are two explanatory variables, $k=2,$ so that $%
y=X_{1}\beta _{1}+X_{2}\beta _{2}+\varepsilon .$ State without proof the
Frisch-Waugh-Lovell theorem for the OLS estimator $\hat{\beta}_{1}.$

(b) Consider two economists estimating the following regression models:
Economist 1 estimates $y=X_{1}\beta _{1}+\varepsilon $ while Economist 2
also includes $X_{2}$ and estimates instead $y=X_{1}\beta _{1}+X_{2}\beta
_{2}+\varepsilon .$ What are the estimators of the two economists for $\beta
_{1}$?

(c) Given that the true data generating process is $y=X_{1}\beta
_{1}+X_{2}\beta _{2}+\varepsilon ,$ show that that the estimator for $\beta
_{1}$ of Economist 2 is unbiased.

(d) Given that the true data generating process is $y=X_{1}\beta
_{1}+X_{2}\beta _{2}+\varepsilon ,$ show that the estimator for $\beta _{1}$
of Economist 1 is biased. Discuss this result in relation to omitted
variable bias.

(e)\ Now suppose the true data generating process is $y=X_{1}\beta
_{1}+\varepsilon ,$ so $\beta _{2}=0.$ Are the two estimators of the two
economists unbiased?

(f) Given that the true $\beta _{2}=0,$ compare the variance of $\hat{\beta}%
_{1}$ for the two economists. What can you conclude about the tradeoff
between bias and variance when deciding which explanatory variables to
include?

\bigskip 

\textbf{Question 3 }

Consider the linear regression model 
\begin{equation*}
y=X\beta +\varepsilon 
\end{equation*}%
where $y$ is an $n\times 1$ vector of dependent variables, $X$ is an $%
n\times k$ matrix of stochastic regressors of rank $k$, $\varepsilon |X\sim 
\mathcal{N}(0,\sigma ^{2}I_{n})$ are random variables, and $\beta $ is a $%
k\times 1$ vector of unknown parameters$.$

(a) Set up the log likelihood function and derive the ML estimators $%
\widehat{\beta }^{ML}$ and $\widehat{\sigma }^{2,ML}.$ Recall the definition
of the density of an $m-$dimensional Gaussian vector $z\sim \mathcal{N}(\mu
,\Omega ):$ 
\begin{equation*}
f\left( z\right) =\left( 2\pi \right) ^{-m/2}\det \left( \Omega \right)
^{-1/2}e^{-\frac{1}{2}\left( z-\mu \right) ^{\prime }\Omega ^{-1}\left(
z-\mu \right) }.
\end{equation*}

(b) Show that $\widehat{\sigma }^{2,ML}$ is biased and propose a bias
corrected estimator for $\sigma ^{2}.$

(c) Show that $\widehat{\sigma }^{2,ML}$ is asymptotically unbiased, and
consistent.

(d) Derive the distribution of $\widehat{\beta }^{ML}|X$ and explain how it
can be used for testing. 

(e) Compute the Fisher information matrix for the $(k+1)\times 1$ vector of
parameters $[\beta ,\sigma ^{2}].$

(f) Are $\widehat{\beta }^{ML}$ and $\widehat{\sigma }^{2,ML}$ efficient?

\end{document}
